\section{Experimental Setup}
\label{sec:experiments}

This section describes the evaluation metrics
(Section~\ref{sec:metrics}), the hardware and software environment
(Section~\ref{sec:env}), and the specific experiments performed
(Section~\ref{sec:exp-design}).

% -------------------------------------------------------------------------
\subsection{Evaluation Metrics}
\label{sec:metrics}
% -------------------------------------------------------------------------

We use the standard COCO detection metrics as implemented by the
Ultralytics evaluation engine:

\begin{itemize}
  \item \textbf{Precision (P)} — the fraction of predicted boxes that
        overlap a ground-truth box with IoU $\ge 0.5$ and the correct class.
  \item \textbf{Recall (R)} — the fraction of ground-truth boxes correctly
        retrieved at the same IoU and class criterion.
  \item \textbf{F1 score} — the harmonic mean of precision and recall,
        computed as $F_1 = 2PR/(P+R)$.
  \item \textbf{mAP$_{50}$} — mean Average Precision at a single IoU
        threshold of 0.50, averaged over classes.
  \item \textbf{mAP$_{50:95}$} — mean Average Precision averaged over
        ten IoU thresholds from 0.50 to 0.95 in steps of 0.05, then
        averaged over classes.  This is the primary COCO metric.
\end{itemize}

All metrics are computed on the \textbf{test split} as the held-out
evaluation set.
The validation split is used only for checkpoint selection (best
mAP$_{50}$ epoch).
Per-class metrics (Precision, Recall, F1, AP$_{50}$, AP$_{50:95}$) are
reported for each of the three classes: \textit{bus}, \textit{car}, and
\textit{van}.

\textbf{TIDE-style error taxonomy.}
In addition to summary AP metrics, we apply a TIDE-inspired error analysis
\cite{bolya2020tide} to decompose test-set errors into six mutually exclusive
categories.
Our implementation in \texttt{stage\_errors} of
\texttt{trace\_experiments.py} follows the assignment rule of
Bolya~et~al.\ (2020) and uses the thresholds defined below.
For each predicted box (sorted by confidence, highest first):

\begin{enumerate}
  \item \textbf{True Positive (TP)} — highest-IoU unmatched ground-truth
        box of the same class with IoU $\ge 0.5$.
  \item \textbf{Duplicate} — same class and IoU $\ge 0.5$, but the
        best matching ground-truth box is already matched by a
        higher-confidence prediction.
  \item \textbf{Classification} — a ground-truth box of a different class
        has IoU $\ge 0.5$; the box is well localised but the class is wrong.
  \item \textbf{Localisation} — same class exists, IoU $\in [0.1, 0.5)$;
        the class is correct but the box does not overlap enough.
  \item \textbf{Both (cls+loc)} — a ground-truth box of a different class
        has IoU $\in [0.1, 0.5)$; both class and location are wrong.
  \item \textbf{Background} — no ground-truth box has IoU $\ge 0.1$;
        the prediction is a spurious detection.
\end{enumerate}

Ground-truth boxes not matched by any prediction are counted as
\textbf{Missed}.
Error counts are stratified by class (\textit{bus}, \textit{car},
\textit{van}, \textit{all}) and by object size
(\textit{small}: area~$<32^2$, \textit{medium}: $32^2$–$96^2$,
\textit{large}: $>96^2$ pixels in the original image).

The confidence threshold for the TIDE analysis is 0.40, matching the
application default (\texttt{DEFAULT\_CONF} in \texttt{app/config.py}).
The foreground IoU threshold is $\tau_{fg} = 0.5$ and the background
threshold is $\tau_{bg} = 0.1$.

\textbf{Latency metrics.}
Inference latency is measured in milliseconds as the wall-clock time for a
single image, averaged over a sample of test images after a warm-up phase of
five images.
We report:
(a)~mean end-to-end latency (ms),
(b)~median end-to-end latency (ms),
(c)~95th-percentile latency (ms), and
(d)~throughput (images per second, FPS) = $1000\,/\,\bar{t}$.
Ultralytics also reports internal stage latencies: pre-process,
inference, and post-process.

% -------------------------------------------------------------------------
\subsection{Hardware and Software Environment}
\label{sec:env}
% -------------------------------------------------------------------------

All experiments reported in this paper were executed on the following
system (Table~\ref{tab:env}):

\begin{table}[htbp]
\centering
\caption{Hardware and software environment (source: stage \texttt{env}
  of \texttt{trace\_experiments.py}).}
\label{tab:env}
\begin{tabular}{ll}
\toprule
Item & Value \\
\midrule
Operating system  & macOS 15.3.1 (arm64) \\
Processor         & Apple M4 Pro (ARM, 14 threads) \\
RAM               & 24.0 GB \\
GPU               & Apple Silicon MPS (no NVIDIA/CUDA) \\
Python            & 3.14.2 \\
PyTorch           & 2.12.0 \\
Ultralytics       & 8.4.53 \\
CUDA              & not available \\
\bottomrule
\end{tabular}
\end{table}

No NVIDIA GPU was available; GPU-accelerated inference used the Apple
Metal Performance Shaders (MPS) backend.
CPU inference was used as the fallback for latency benchmarking, and all
training runs were executed on CPU (the \texttt{train} stage was not
executed in the current run set; see \texttt{run\_args.json}).

The experiment harness automatically selects the best available device:
CUDA~GPU if available, then MPS (Apple Silicon), then CPU.
This logic is implemented in \texttt{auto\_device()} in
\texttt{trace\_experiments.py}.

% -------------------------------------------------------------------------
\subsection{Experiment Design}
\label{sec:exp-design}
% -------------------------------------------------------------------------

The full experiment harness (\texttt{trace\_experiments.py}) is organised
into seven stages, each producing one or more output tables and figures:

\begin{description}
  \item[\texttt{audit}]
    Label-format audit and production of a cleaned dataset copy
    (Tables~T2, B1; Section~\ref{sec:audit}).

  \item[\texttt{env}]
    Hardware and software environment snapshot (Table~T5).

  \item[\texttt{eval}]
    Evaluation of the deployed \texttt{best.pt} checkpoint on the validation
    and test splits, for both original and cleaned labels.
    Produces overall metrics (Table~T6), per-class metrics (Table~T7),
    confusion-matrix plots (Figure~F8), and PR~curves (Figure~F9).

  \item[\texttt{latency}]
    Inference-latency and throughput measurement on CPU
    (and MPS when available) over up to 200 test images, plus model-size
    statistics (Table~T10).

  \item[\texttt{errors}]
    TIDE-style error taxonomy on the test set (Table~T11, T11b) and
    qualitative success/failure figures (Figure~F10).

  \item[\texttt{train}] \emph{Not executed in the current run.}
    Re-trains YOLO26m, YOLOv8m, YOLO11m, and RT-DETR-l from COCO
    pre-trained weights on the cleaned dataset, over three seeds, with the
    photometric-only augmentation policy.
    Produces Table~T8 (baseline comparison) and Table~T9 (augmentation
    ablation for YOLO26m).
    Results are marked [DATA DIBUTUHKAN] throughout this paper.

  \item[\texttt{report}]
    Collects all tables and figure paths into
    \texttt{results/summary.md} and archives outputs to
    \texttt{results.zip}.
\end{description}

The stages executed in the current run were:
\texttt{audit}, \texttt{env}, \texttt{eval}, \texttt{latency},
\texttt{errors}, and \texttt{report}.
This covers the evaluation of the already-trained deployed checkpoint
but excludes new training.
The command used was:

\begin{verbatim}
python3 trace_experiments.py \
  --data-root ./data \
  --weights app/models/best.pt \
  --out ./results \
  --device mps \
  --stages audit env eval latency errors report
\end{verbatim}

Completed stages are skipped automatically on re-runs; incremental
execution allows adding the \texttt{train} stage at a later date without
repeating earlier stages.

\textbf{Reproducibility.}
All random seeds are fixed (\texttt{deterministic=True}) in
\texttt{BASE\_TRAIN}.
Training seeds 0, 1, and 2 are planned for the baseline and ablation
experiments.
The cleaned dataset is reproduced deterministically by the \texttt{audit}
stage from the original Roboflow export.
Model weights, the experiment harness, the CI workflow, and the test suite
are stored in the public GitHub repository
\texttt{indri007/TRACE}\footnote{Dataset files and model weights are
excluded from version control via \texttt{.gitignore}.}.
